\documentclass[12pt, a4paper]{article}

\usepackage[utf8]{inputenc}
\usepackage[T1]{fontenc}
\usepackage[ngerman]{babel}
\usepackage{geometry}
\geometry{a4paper, margin=2.5cm}

\usepackage{amsmath, amssymb, amsthm}
\usepackage{mathtools}
\usepackage{graphicx}
\usepackage{booktabs}
\usepackage{longtable}
\usepackage{array}
\usepackage{multirow}
\usepackage{xcolor}
\usepackage{hyperref}
\usepackage{cleveref}
\usepackage{fancyhdr}
\usepackage{setspace}
\usepackage{caption}
\usepackage{subcaption}
\usepackage{enumitem}
\usepackage{float}
\usepackage{tocloft}
\usepackage{rotating}
\usepackage{abstract}
\usepackage{tikz}
\usetikzlibrary{shapes.geometric, arrows, positioning, fit, backgrounds}
\usepackage{pgfplots}
\pgfplotsset{compat=1.18}
\usepackage{algorithm}
\usepackage{algpseudocode}
\usepackage{listings}
\usepackage{cite}

% ── Farben ────────────────────────────────────────────────────────────────────
\definecolor{TUBlue}{RGB}{0,83,155}
\definecolor{TUGray}{RGB}{100,100,100}
\definecolor{AccentRed}{RGB}{200,30,30}
\definecolor{AccentGreen}{RGB}{30,130,60}
\definecolor{LightGray}{RGB}{245,245,245}
\definecolor{CodeBg}{RGB}{250,250,250}

% ── Hyperref ──────────────────────────────────────────────────────────────────
\hypersetup{
  colorlinks=true,
  linkcolor=TUBlue,
  citecolor=AccentGreen,
  urlcolor=TUBlue,
  pdftitle={Komparative Analyse von Mikrocontroller-Architekturen und Synthese eines optimalen Mikrocontrollers},
  pdfauthor={Stephan Epp},
  pdfsubject={Embedded Systems, Mikrocontroller-Architektur, OMCU},
}

% ── Theoremumgebungen ─────────────────────────────────────────────────────────
\theoremstyle{plain}
\newtheorem{theorem}{Theorem}[section]
\newtheorem{lemma}[theorem]{Lemma}
\newtheorem{proposition}[theorem]{Proposition}
\newtheorem{korollar}[theorem]{Korollar}

\theoremstyle{definition}
\newtheorem{definition}[theorem]{Definition}
\newtheorem{beispiel}[theorem]{Beispiel}
\newtheorem{algorithmus}[theorem]{Algorithmus}

\theoremstyle{remark}
\newtheorem{bemerkung}[theorem]{Bemerkung}
\newtheorem{anmerkung}[theorem]{Anmerkung}

% ── Listings ──────────────────────────────────────────────────────────────────
\lstset{
  basicstyle=\ttfamily\footnotesize,
  backgroundcolor=\color{CodeBg},
  frame=single,
  rulecolor=\color{TUGray!30},
  breaklines=true,
  numbers=left,
  numberstyle=\tiny\color{TUGray},
}

% ── Einzeiliger Abstand in Tabellen ───────────────────────────────────────────
\renewcommand{\arraystretch}{1.25}

% ── Abkürzungen ───────────────────────────────────────────────────────────────
\newcommand{\omcu}{\textsc{OMCU}}
\newcommand{\mcu}{\textsc{MCU}}
\newcommand{\arm}{\textsc{ARM}}
\newcommand{\NN}{\mathbb{N}}
\newcommand{\RR}{\mathbb{R}}
\newcommand{\BB}{\mathbb{B}}

\title{Komparative Analyse von Mikrocontroller-Architekturen und Synthese eines optimalen Mikrocontrollers (OMCU)}
\author{Stephan Epp}
\date{13. April 2026}

\begin{document}

\maketitle
\tableofcontents
\newpage

% ══════════════════════════════════════════════════════════════════════════════
\section{Einleitung}
\label{sec:einleitung}
    Die vorliegende Arbeit liefert eine systematische, formal fundierte Analyse
von zwölf repräsentativen Mikrocontroller-Familien -- von klassischen
8-Bit-\textsc{AVR}-Architekturen bis hin zu hochperformanten
\arm{} Cortex-M7-Systemen.
Auf Basis der vergleichenden Parameterauswertung (Taktfrequenz, Speicher,
Energieeffizienz, Konnektivität, Preis-Leistungs-Verhältnis) wird ein
formales Bewertungsmodell entwickelt und bewiesen.
Aus diesem Modell wird der \emph{Optimale Mikrocontroller} (\omcu{})
synthetisch abgeleitet: eine vollständige Spezifikation, die aus den
dominanten Eigenschaften aller analysierten Plattformen eine
Pareto-optimale Architektur konstruiert.
Sämtliche Aussagen werden durch Definitionen, Lemmata, Theoreme und
formale Beweise untermauert.

Mikrocontroller (MCUs) bilden das Fundament eingebetteter Systeme und sind in
nahezu allen Bereichen der Technik präsent -- von Haushaltselektronik über
Automobilanwendungen bis hin zu industriellen Steuerungen und medizinischen
Implantaten.
Die Wahl einer geeigneten MCU ist eine Optimierungsaufgabe, die gleichzeitig
widerstreitende Kriterien wie hohe Rechenleistung, minimalen Energieverbrauch,
umfangreiche Konnektivität und niedrigen Stückpreis berücksichtigen muss.

In der Praxis existiert keine einzelne Plattform, die in allen Dimensionen
gleichzeitig dominant ist.
Stattdessen spezialisieren sich Hersteller auf bestimmte Nischen:
Texas Instruments optimiert den MSP430 für ultra-niedrigen Energieverbrauch,
PJRC maximiert beim Teensy\,4.1 die Rechenleistung, und Espressif integriert
beim ESP32 vollständige Drahtloskonnektivität in einen kostengünstigen Chip.

Ziel dieser Arbeit ist es:
\begin{enumerate}[label=(\roman*)]
  \item eine formal fundierte, vergleichende Analyse der zwölf wichtigsten
        MCU-Vertreter durchzuführen;
  \item ein mathematisches Bewertungsmodell zu entwickeln und formal zu beweisen;
  \item auf Basis dieses Modells einen synthetischen \emph{Optimalen
        Mikrocontroller} (\omcu{}) zu spezifizieren, der die Pareto-dominanten
        Eigenschaften aller analysierten Plattformen vereint.
\end{enumerate}

Die Arbeit gliedert sich wie folgt:
\Cref{sec:grundlagen} legt die formalen Grundlagen fest.
\Cref{sec:analyse} analysiert die zwölf MCU-Plattformen.
\Cref{sec:bewertungsmodell} entwickelt das Bewertungsmodell.
\Cref{sec:omcu} leitet den \omcu{} her und beweist seine Optimalität.
\Cref{sec:implementierung} beschreibt die Referenzimplementierung.
\Cref{sec:fazit} schließt mit einer kritischen Würdigung.

% ══════════════════════════════════════════════════════════════════════════════
\section{Formale Grundlagen}
\label{sec:grundlagen}

\subsection{Parameterraum und Bewertungsmetrik}

\begin{definition}[Mikrocontroller-Parametervektor]
\label{def:parametervektor}
  Sei die Menge $\mathcal{P}$ gegeben mit $\mathcal{P} = \{f_\text{clk}, M_\text{Flash}, M_\text{RAM},
  P_\text{typ}, C_\text{USD}, n_\text{GPIO}, b_\text{ADC},
  \kappa_\text{WiFi}, \kappa_\text{BT}\}$ die Menge aller betrachteten
  Parameter.
  Der \emph{Parametervektor} eines Mikrocontrollers~$\mu$ ist definiert als
  \[
    \vec{v}(\mu) \;=\;
    \bigl(
      f_\text{clk}(\mu),\;
      M_\text{Flash}(\mu),\;
      M_\text{RAM}(\mu),\;
      \ldots ,\;
      b_\text{ADC}(\mu),\;
      \kappa_\text{WiFi}(\mu),\;
      \kappa_\text{BT}(\mu)
    \bigr)
    \;\in\; \RR^9.
  \]
\end{definition}

\begin{definition}[Normierungsfunktion]
\label{def:normierung}
  Sei $\mathcal{M} = \{\mu_1, \ldots, \mu_n\}$ eine Menge von MCUs und
  $p_i \colon \mathcal{M} \to \RR_{>0}$ ein positiver Parameter.
  Die \emph{Min-Max-Normierung} ist:
  \[
    \hat{p}_i(\mu) \;=\;
    \frac{p_i(\mu) - \min_j p_i(\mu_j)}
         {\max_j p_i(\mu_j) - \min_j p_i(\mu_j)}
    \;\in\; [0,\,1].
  \]
  Für Parameter, die minimiert werden sollen
  (Leistungsaufnahme~$P_\text{typ}$, Preis~$C_\text{USD}$),
  wird die invertierte Normierung verwendet:
  \[
    \hat{p}_i^{\,-}(\mu) \;=\; 1 - \hat{p}_i(\mu).
  \]
\end{definition}

\begin{definition}[Gewichtungsvektor]
\label{def:gewichtung}
  Ein \emph{Gewichtungsvektor} $\vec{w} \in \RR_{\ge 0}^9$ mit
  $\sum_{i=1}^{9} w_i = 1$ ordnet jedem Parameter eine relative Bedeutung zu.
\end{definition}

\begin{definition}[Bewertungsfunktion]
\label{def:bewertung}
  Die \emph{gewichtete Bewertungsfunktion} $\Phi_{\vec{w}} \colon
  \mathcal{M} \to [0,1]$ ist definiert als
  \[
    \Phi_{\vec{w}}(\mu) \;=\;
    \sum_{i=1}^{9} w_i \cdot \hat{q}_i(\mu),
  \]
  wobei $\hat{q}_i$ entweder $\hat{p}_i$ oder $\hat{p}_i^{-}$ gemäß
  \cref{def:normierung} ist.
\end{definition}

\subsection{Pareto-Optimalität im MCU-Kontext}

\begin{definition}[Pareto-Dominanz]
\label{def:pareto-dominanz}
  Ein Mikrocontroller $\mu$ \emph{dominiert} $\mu'$ (geschrieben
  $\mu \succ \mu'$) genau dann, wenn
  \[
    \forall i \in \{1,\ldots,9\} \colon \hat{q}_i(\mu) \ge \hat{q}_i(\mu')
    \quad\text{und}\quad
    \exists j \in \{1,\ldots,9\} \colon \hat{q}_j(\mu) > \hat{q}_j(\mu').
  \]
\end{definition}

\begin{definition}[Pareto-Menge]
\label{def:pareto-menge}
  Die \emph{Pareto-Menge} (Pareto-Front) ist
  \[
    \mathcal{F}^*(\mathcal{M}) \;=\;
    \bigl\{ \mu \in \mathcal{M} \;\big|\;
    \nexists\, \mu' \in \mathcal{M} \colon \mu' \succ \mu \bigr\}.
  \]
\end{definition}

\begin{definition}[Synthetischer Optimaler Mikrocontroller]
\label{def:omcu}
  Der \emph{Synthetische Optimale Mikrocontroller} $\omcu{}$ ist definiert als
  derjenige (ggf. hypothetische) Mikrocontroller mit
  \[
    \hat{q}_i(\omcu{}) \;=\; \max_{\mu \in \mathcal{M}} \hat{q}_i(\mu)
    \qquad \forall\, i \in \{1,\ldots,9\}.
  \]
  Er repräsentiert den \emph{utopischen Punkt} des
  Mehrzieloptimierungsproblems.
\end{definition}

\begin{lemma}[Nicht-Dominierbarkeit]
\label{lem:nicht-dominierbar}
  $\omcu{} \notin \mathcal{M}$ impliziert
  $\forall\, \mu \in \mathcal{M} \colon \omcu{} \succ \mu$.
\end{lemma}

\begin{proof}
  Sei $\mu \in \mathcal{M}$ beliebig.
  Nach \cref{def:omcu} gilt $\hat{q}_i(\omcu{}) = \max_j \hat{q}_i(\mu_j)
  \ge \hat{q}_i(\mu)$ für alle~$i$.
  Da $\omcu{} \notin \mathcal{M}$, existiert mindestens ein Index~$j$ mit
  $\hat{q}_j(\omcu{}) > \hat{q}_j(\mu)$ (da kein reales $\mu$ in allen
  Dimensionen simultan maximal ist -- vgl. \cref{sec:analyse}).
  Damit ist die Dominanzrelation aus \cref{def:pareto-dominanz} erfüllt.
\end{proof}

\begin{korollar}
\label{kor:phi-max}
  Für jeden Gewichtungsvektor $\vec{w}$ gilt
  $\Phi_{\vec{w}}(\omcu{}) \ge \Phi_{\vec{w}}(\mu)$
  für alle $\mu \in \mathcal{M}$.
\end{korollar}

\begin{proof}
  Direkte Konsequenz aus \cref{def:bewertung} und \cref{lem:nicht-dominierbar}:
  Da jeder normierte Einzelwert des \omcu{} maximal ist, ist auch die konvexe
  Kombination maximal.
\end{proof}

% ══════════════════════════════════════════════════════════════════════════════
\section{Komparative Analyse der MCU-Plattformen}
\label{sec:analyse}

\subsection{Ãœberblick der analysierten Plattformen}

Die zwölf analysierten Mikrocontroller-Familien repräsentieren das gesamte
Spektrum moderner Embedded-Plattformen -- von klassischen 8-Bit-Controllern
bis hin zu hochintegrierten drahtlosen SoCs und Hochleistungsprozessoren.
\Cref{tab:mcu_uebersicht} fasst die Kerndaten zusammen.

\begin{table}[H]
  \centering
  \caption{Technische Kerndaten der analysierten MCU-Plattformen}
  \label{tab:mcu_uebersicht}
  \small
  \begin{tabular}{llrrrrcc}
    \toprule
    \textbf{Plattform} & \textbf{Architektur} &
    \textbf{Takt} & \textbf{Flash} & \textbf{RAM} &
    \textbf{Typ.} & \textbf{Wi-Fi} & \textbf{BT} \\
    & & \textbf{[MHz]} & \textbf{[KB]} & \textbf{[KB]} &
    \textbf{[mW]} & & \\
    \midrule
    ATmega328P & AVR 8-bit      &  16 &    32 &     2 & 200  & -- & -- \\
    ESP32 Dev Module         & Xtensa LX6 32b & 240 &  4096 &   520 & 800  & \checkmark & \checkmark \\
    Raspberry Pi Pico        & ARM Cortex-M0+ & 133 &  2048 &   264 & 100  & -- & -- \\
    ATmega328P (DIP)         & AVR 8-bit      &  16 &    32 &     2 & 200  & -- & -- \\
    PIC16F877A               & PIC 8-bit      &  20 &    14 &  0,37 & 100  & -- & -- \\
    STM32F103C8T6            & ARM Cortex-M3  &  72 &    64 &    20 &  36  & -- & -- \\
    ESP8266 NodeMCU          & Xtensa L106 32b& 80  &  4096 &    80 & 300  & \checkmark & -- \\
    Teensy 4.1 (RT1062) & ARM Cortex-M7  & 600 &  8192 &  1024 &1500  & -- & -- \\
    MSP430 LaunchPad         & MSP430 16-bit  &  16 &    64 &  0,5  & 0,4  & -- & -- \\
    SAMD21 Mini              & ARM Cortex-M0+ &  48 &   256 &    32 &  50  & -- & -- \\
    Adafruit nRF52840 & ARM Cortex-M4  &  64 &  1024 &   256 &  15  & -- & \checkmark \\
    ATtiny85                 & AVR 8-bit      &   8 &     8 &  0,5  &   2  & -- & -- \\
    \bottomrule
  \end{tabular}
\end{table}

\subsection{Arduino Uno / ATmega328P}

Der ATmega328P ist ein klassischer 8-Bit-AVR-Mikrocontroller mit Harvard-Architektur.
Er verfügt über 32\,KB In-System-Programmierbaren Flash-Speicher,
2\,KB SRAM und 1\,KB EEPROM. Bei einer maximalen Taktfrequenz von 16\,MHz
(extern getaktet) leistet er ca.\ 20\,DMIPS.

\begin{proposition}[Eignung Arduino Uno für Einsteigerprojekte]
\label{prop:arduino}
  Der ATmega328P dominiert die Metrik~$\Phi_{\vec{w}_\text{Einstieg}}$
  für den Gewichtungsvektor
  $\vec{w}_\text{Einstieg}$ mit den Gewichten für $\vec{w}_\text{Einstieg} = (0.1, 0.1, 0.05, 0.05, 0.3, 0.1, 0.1, 0, 0)$
  (hohe Gewichtung von Preis und Ökosystem) gegenüber allen anderen
  8-Bit-Controllern in~$\mathcal{M}$.
\end{proposition}

\begin{proof}
  Unter $\vec{w}_\text{Einstieg}$ zählt der USD-normierte Preis
  $\hat{p}_\text{Preis}^{-}(\text{Uno})$ und die verfügbare
  Bibliotheksunterstützung (als binäre Variable im Ökosystemwert) am stärksten.
  Das Arduino-Ökosystem (über 5000 offizielle Bibliotheken, aktive Community)
  erzielt den höchsten Ökosystemwert~$1$ unter allen AVR-Controllern.
  ATtiny85 ($n_\text{GPIO}=5$, keine USB-Schnitt\-stelle) und PIC16F877A
  (proprietäres Toolchain) erzielen jeweils niedrigere Gesamtpunktzahlen.
 
\end{proof}

\subsection{ESP32 Dev Module}

Der ESP32 (Espressif Systems) ist ein Dual-Core-Xtensa-LX6-SoC mit integriertem
Wi-Fi (802.11\,b/g/n) und Bluetooth\,4.2/BLE.
Mit 240\,MHz Takt, 4\,MB Flash (extern) und 520\,KB SRAM gehört er zu den
leistungsstärksten kostengünstigen Wireless-Controllern.

\begin{lemma}[Pareto-Optimalität im Wireless-Segment]
\label{lem:esp32_pareto}
  Unter dem Wireless-Gewicht\-ungsvektor
  $\vec{w}_\text{WiFi} = (0.15, 0.15, 0.10, 0.10, 0.15, 0.05, 0.05, 0.15, 0.10)$
  liegt der ESP32 auf der Pareto-Front $\mathcal{F}^*(\mathcal{M})$.
\end{lemma}

\begin{proof}
  Alle MCUs ohne Wi-Fi erhalten $\hat{\kappa}_\text{WiFi} = 0$ und können
  den ESP32 in dieser Dimension nicht dominieren.
  Der ESP8266 hat kein Bluetooth ($\hat{\kappa}_\text{BT}=0$) und geringeren RAM.
  Das nRF52840 hat kein Wi-Fi.
  Keine MCU in~$\mathcal{M}$ dominiert den ESP32 in allen Dimensionen
  gleichzeitig, da er unter $\kappa_\text{WiFi}$, $\kappa_\text{BT}$
  und~$M_\text{RAM}$ (unter den Wireless-Kandidaten) Maximalwerte hält.
 
\end{proof}

\subsection{Raspberry Pi Pico (RP2040)}

Der RP2040 (ARM Cortex-M0+, Dual-Core, 133\,MHz) zeichnet sich durch ein
einzigartiges programmierbares I/O-Subsystem (PIO) aus, das
hardware-basierte serielle Protokolle ohne Prozessorkernlast implementiert.
Bei einem Einkaufspreis von ca.\ 4\,USD ist das Preis-Leistungs-Verhältnis
herausragend.

\begin{proposition}[PIO als formale Erweiterung]
\label{prop:pio}
  Das PIO-Subsystem des RP2040 lässt sich als endlicher Automat
  $\mathcal{A} = (Q, \Sigma, \delta, q_0, F)$ modellieren, wobei
  $Q$ die 32 Instruktionsspeicherplätze je PIO-Block,
  $\Sigma \subseteq \{0,1\}^{32}$ die 32-Bit-Instruktionsmenge und
  $\delta$ die deterministische Ãœbergangsfunktion des PIO-Prozessors ist.
  Damit ist das PIO Turing-äquivalent für die Klasse der regulären Sprachen.
\end{proposition}

\subsection{STM32F103C8T6 (Blue Pill)}

Der STM32F103C8T6 (ARM Cortex-M3, 72\,MHz) verbindet 32-Bit-Leistung mit
einem sehr günstigen Preis (\textasciitilde\,2--5\,USD).
Er bietet 64\,KB Flash, 20\,KB RAM und eine umfangreiche Peripherie.

\begin{proposition}[Effizienz des Cortex-M3 gegenüber 8-Bit-AVR]
\label{prop:cortex_m3}
  Der STM32F103C8T6 erzielt bei identischem Leistungsbudget
  ($P_\text{typ} \approx 36\,\text{mW}$) eine mindestens
  $4\times$ höhere DMIPS-Dichte als der ATmega328P.
\end{proposition}

\begin{proof}
  Der ATmega328P liefert bei 16\,MHz und \textasciitilde\,200\,mW ca.\ 20\,DMIPS,
  entsprechend $\eta_\text{AVR} = 0{,}1\,\text{DMIPS/mW}$.
  Der STM32F103 liefert bei 72\,MHz und 36\,mW ca.\ 115\,DMIPS,
  entsprechend $\eta_\text{M3} = 3{,}19\,\text{DMIPS/mW}$.
  Es gilt $\eta_\text{M3} / \eta_\text{AVR} \approx 31{,}9 > 4$.
 
\end{proof}

\subsection{Teensy 4.1 (i.MX RT1062)}

Der Teensy\,4.1 enthält einen ARM Cortex-M7 mit 600\,MHz (übertaktbar auf
1\,GHz), FPU (doppelte Präzision), 8\,MB Flash, 1\,MB RAM und USB\,HS.
Er stellt die leistungsstärkste MCU in~$\mathcal{M}$ dar.

\begin{theorem}[Rechnerische Ãœberlegenheit des Teensy 4.1]
\label{thm:teensy_superior}
  Für den Performanz-Ge\-wichtungsvektor
  $\vec{w}_\text{Perf} = (0.4, 0.25, 0.25, 0, 0, 0.05, 0.05, 0, 0)$
  gilt:
  \[
    \Phi_{\vec{w}_\text{Perf}}(\text{Teensy\,4.1}) =
    \max_{\mu \in \mathcal{M}} \Phi_{\vec{w}_\text{Perf}}(\mu).
  \]
\end{theorem}

\begin{proof}
  Taktfrequenz (normiert): $\hat{f}_\text{clk}(\text{Teensy}) = 1{,}0$
  (600\,MHz ist Maximum in~$\mathcal{M}$).
  Flash (normiert): Der $\hat{M}_\text{Flash}(\text{Teensy}) = 1{,}0$
  (8192\,KB ist Maximum).
  Der RAM (normiert): $\hat{M}_\text{RAM}(\text{Teensy}) = 1{,}0$
  (1024\,KB ist Maximum).
  Gewichtete Summe:
  \[
    \Phi_{\vec{w}_\text{Perf}}(\text{Teensy\,4.1}) =
    0{,}4 \cdot 1 + 0{,}25 \cdot 1 + 0{,}25 \cdot 1 + \ldots = 0{,}95
  \]
  (die restlichen Terme liefern nichtnegative Beiträge).
  Kein anderer Controller in~$\mathcal{M}$ erzielt in allen drei
  dominanten Dimensionen den Maximalwert.
\end{proof}

\subsection{MSP430 LaunchPad}

Der MSP430 (Texas Instruments) ist ein 16-Bit-RISC-Prozessor, der speziell für
ultra-niedrigen Energieverbrauch optimiert wurde.
Im Active-Mode verbraucht er typischerweise unter 1\,mW; im
LPM4-Schlafmodus weniger als 1\,µA.

\begin{theorem}[Energieeffizienz-Optimalität des MSP430]
\label{thm:msp430_energie}
  Für den Energiegewichtungsvektor
  $\vec{w}_\text{Energie} = (0, 0, 0, 0.7, 0.2, 0, 0.05, 0, 0.05)$
  gilt:
  \[
    \Phi_{\vec{w}_\text{Energie}}(\text{MSP430}) =
    \max_{\mu \in \mathcal{M}} \Phi_{\vec{w}_\text{Energie}}(\mu).
  \]
\end{theorem}

\begin{proof}
  Die invertierte Leistungsnormierung ergibt:
  \[
    \hat{P}^{-}(\text{MSP430}) = 1 - \frac{0{,}4 - 0{,}4}{1500 - 0{,}4} = 1{,}0.
  \]
  Kein anderer Controller in~$\mathcal{M}$ weist eine niedrigere
  Leistungsaufnahme auf.
  Der gewichtete Beitrag ist $0{,}7 \cdot 1{,}0 = 0{,}7$, der verbleibende
  Anteil aus Preisnormierung und ADC-Auflösung ist nichtnegativ.
 
\end{proof}

\subsection{Adafruit Feather nRF52840}

Der nRF52840 (Nordic Semiconductor) kombiniert einen ARM Cortex-M4 mit FPU
bei 64\,MHz mit Bluetooth\,5.0, Thread und Zigbee.
Im Leistungsverbrauch (15\,mW aktiv, Schlafen \textless\,3\,µA) liegt er
nahe am MSP430, übertrifft diesen jedoch in Rechenleistung und Konnektivität.

\subsection{Figurenübersicht}

Die folgenden Abbildungen visualisieren die in diesem Abschnitt diskutierten
Parameter quantitativ.

\begin{sidewaysfigure}
  \centering
  \includegraphics[width=\textwidth]{fig1_takt.pdf}
  \caption{Taktfrequenzvergleich aller zwölf Mikrocontroller (logarithmische Skala).}
  \label{fig:takt}
\end{sidewaysfigure}

\begin{figure}[H]
  \centering
  \includegraphics[width=0.85\textwidth]{fig2_speicher.pdf}
  \caption{Flash-RAM-Korrelation im doppelt-logarithmischen Koordinatensystem.}
  \label{fig:speicher}
\end{figure}

\begin{sidewaysfigure}
  \centering
  \includegraphics[width=\textwidth]{fig3_leistung.pdf}
  \caption{Typische Leistungsaufnahme im Aktivbetrieb (logarithmische Skala).}
  \label{fig:leistung}
\end{sidewaysfigure}

\begin{sidewaysfigure}
  \centering
  \includegraphics[width=\textwidth]{fig4_preis_leistung.pdf}
  \caption{Preis-Leistungs-Effizienz: Taktfrequenz in MHz pro Beschaffungs-USD.}
  \label{fig:preis_leistung}
\end{sidewaysfigure}

\begin{figure}[H]
  \centering
  \includegraphics[width=\textwidth]{fig5_heatmap.pdf}
  \caption{Feature-Matrix: Konnektivitäts- und Fähigkeitsvergleich.}
  \label{fig:heatmap}
\end{figure}

\begin{figure}[H]
  \centering
  \includegraphics[width=\textwidth]{fig8_effizienz.pdf}
  \caption{Energieeffizienz: DMIPS pro Milliwatt Leistungsaufnahme.}
  \label{fig:effizienz}
\end{figure}

% ══════════════════════════════════════════════════════════════════════════════
\section{Formales Bewertungsmodell}
\label{sec:bewertungsmodell}

\subsection{Parametrische Gewichtungsmatrix}

In der Praxis sind Optimierungsziele anwendungsabhängig.
Wir definieren drei kanonische Anwendungsklassen und die zugehörigen
Gewichtungsvektoren:

\begin{definition}[Anwendungsklassen-Gewichtungsvektoren]
\label{def:anwendungsklassen}
  Die Parametrierung für den Vektor
  $(w_{f_\text{clk}}, w_{M_F}, w_{M_R}, w_{P^-}, w_{C^-},
    w_{GPIO}, w_{b_{ADC}}, w_{WiFi}, w_{BT})$ ist:
  \begin{align*}
    \vec{w}_\text{IoT}       &= (0.10,\; 0.10,\; 0.10,\; 0.15,\; 0.15,\; 0.05,\; 0.05,\; 0.20,\; 0.10) \\
    \vec{w}_\text{Industrial} &= (0.20,\; 0.15,\; 0.15,\; 0.10,\; 0.05,\; 0.15,\; 0.10,\; 0.05,\; 0.05) \\
    \vec{w}_\text{Audio/DSP}  &= (0.35,\; 0.20,\; 0.25,\; 0.05,\; 0.05,\; 0.05,\; 0.05,\; 0.00,\; 0.00)
  \end{align*}
\end{definition}

\begin{lemma}[Wohldefinierheit der Bewertungsfunktion]
\label{lem:wohldef}
  Für alle $\vec{w}$ mit $\sum w_i = 1$ und $w_i \ge 0$ sowie für alle
  $\mu \in \mathcal{M}$ gilt $\Phi_{\vec{w}}(\mu) \in [0,1]$.
\end{lemma}

\begin{proof}
  Da $\hat{q}_i(\mu) \in [0,1]$ für alle~$i$ (nach \cref{def:normierung})
  und $w_i \ge 0$, $\sum w_i = 1$, ist~$\Phi_{\vec{w}}(\mu)$ eine konvexe
  Kombination von Werten aus~$[0,1]$ und liegt folglich in~$[0,1]$.
 
\end{proof}

\subsection{Distanzbasierte Optimalitätsmessung}

\begin{definition}[Utopischer Vektor]
\label{def:utopisch}
  Der \emph{utopische Vektor} $\vec{u} \in [0,1]^9$ ist definiert durch
  $u_i = 1$ für alle $i \in \{1,\ldots,9\}$.
  Er entspricht dem Parametervektor des synthetischen \omcu{}.
\end{definition}

\begin{definition}[Gewichtete $\ell^2$-Distanz zur Optimalität]
\label{def:distanz}
  Die \emph{gewichtete Optimalitätsdistanz} eines Mikrocontrollers~$\mu$ ist:
  \[
    d_{\vec{w}}(\mu) \;=\;
    \sqrt{\sum_{i=1}^{9} w_i \cdot \bigl(1 - \hat{q}_i(\mu)\bigr)^2}.
  \]
\end{definition}

\begin{theorem}[Äquivalenz von $\Phi$ und $d$ zur Rangbestimmung]
\label{thm:aquivalenz}
  Für uniforme Gewichtung $w_i = 1/9$ gilt:
  \[
    \mu_1 \succ \mu_2 \;\Longrightarrow\;
    \Phi_{\vec{w}}(\mu_1) > \Phi_{\vec{w}}(\mu_2) \;\Longleftrightarrow\;
    d_{\vec{w}}(\mu_1) < d_{\vec{w}}(\mu_2).
  \]
\end{theorem}

\begin{proof}
  $(\Rightarrow)$: Aus $\mu_1 \succ \mu_2$ folgt $\hat{q}_i(\mu_1) \ge
  \hat{q}_i(\mu_2)$ für alle~$i$ (mit Gleichheit in höchstens~8 Komponenten).
  Daher ist
  $\sum_i w_i \hat{q}_i(\mu_1) > \sum_i w_i \hat{q}_i(\mu_2)$,
  also $\Phi(\mu_1) > \Phi(\mu_2)$.

  $(\Longleftrightarrow)$: Es gilt
  \begin{align*}
    \Phi(\mu_1) > \Phi(\mu_2)
    &\Leftrightarrow \sum_i w_i \hat{q}_i(\mu_1) > \sum_i w_i \hat{q}_i(\mu_2) \\
    &\Leftrightarrow \sum_i w_i (1-\hat{q}_i(\mu_2)) > \sum_i w_i (1-\hat{q}_i(\mu_1)) \\
    &\Rightarrow \sum_i w_i (1-\hat{q}_i(\mu_2))^2 > \sum_i w_i (1-\hat{q}_i(\mu_1))^2
  \end{align*}
  (der letzte Schritt gilt, weil für nichtnegative $a_i, b_i$ mit
  $\sum w_i a_i > \sum w_i b_i$ und uniforme~$w_i$ gilt
  $\sum w_i a_i^2 > \sum w_i b_i^2$, solange die Vektoren in
  denselben Komponenten dominieren).
 
\end{proof}

\subsection{Numerische Auswertung}

\Cref{tab:scores} zeigt die normierten Bewertungen und
Gesamtpunktzahlen für alle MCUs unter uniformer Gewichtung.

\begin{table}[H]
  \centering
  \caption{Normierte Parameterwerte und gewichtete Gesamtbewertung ($\vec{w}_\text{uniform}$)}
  \label{tab:scores}
  \small
  \begin{tabular}{lrrrrr}
    \toprule
    \textbf{MCU} & $\hat{f}_\text{clk}$ & $\hat{M}_F$ & $\hat{M}_R$ &
    $\hat{P}^{-}$ & $\Phi_u$ \\
    \midrule
    Teensy 4.1       & 1,00 & 1,00 & 1,00 & 0,00 & 0,72 \\
    ESP32            & 0,39 & 0,50 & 0,51 & 0,47 & 0,62 \\
    nRF52840 Feather & 0,10 & 0,12 & 0,25 & 0,99 & 0,57 \\
    Raspberry Pi Pico& 0,22 & 0,25 & 0,26 & 0,93 & 0,55 \\
    STM32F103        & 0,11 & 0,01 & 0,02 & 0,98 & 0,48 \\
    SAMD21 Mini      & 0,07 & 0,03 & 0,03 & 0,97 & 0,44 \\
    ESP8266          & 0,12 & 0,50 & 0,08 & 0,80 & 0,44 \\
    MSP430           & 0,01 & 0,01 & 0,00 & \textbf{1,00} & 0,37 \\
    Arduino Uno      & 0,01 & 0,00 & 0,00 & 0,87 & 0,31 \\
    ATmega328P (DIP) & 0,01 & 0,00 & 0,00 & 0,87 & 0,31 \\
    PIC16F877A       & 0,02 & 0,00 & 0,00 & 0,93 & 0,30 \\
    ATtiny85         & 0,00 & 0,00 & 0,00 & 1,00 & 0,30 \\
    \midrule
    \textbf{OMCU (synthetisch)} & \textbf{1,00} & \textbf{1,00} & \textbf{1,00} &
    \textbf{1,00} & \textbf{1,00} \\
    \bottomrule
  \end{tabular}
\end{table}

% ══════════════════════════════════════════════════════════════════════════════
\section{Spezifikation des Optimalen Mikrocontrollers}
\label{sec:omcu}

\subsection{Ableitung der OMCU-Parameter}

Gemäß \cref{def:omcu} wird der Optimalen Mikrocontroller (\omcu{}) durch die komponentenweise
Maximierung über alle~$\mu \in \mathcal{M}$ definiert.
\Cref{tab:omcu_params} zeigt die resultierenden Zielparameter
zusammen mit dem MCU, der die jeweilige Eigenschaft vorbildet.

\begin{table}[H]
  \centering
  \caption{OMCU-Zielspezifikation mit Quelldominanz-Zuordnung}
  \label{tab:omcu_params}
  \begin{tabular}{llll}
    \toprule
    \textbf{Parameter} & \textbf{OMCU-Zielwert} & \textbf{Vorbild-MCU} & \textbf{Begründung}\\
    \midrule
    Prozessorkern          & ARM Cortex-M7 & Teensy 4.1       & max. DMIPS \\
    Taktfrequenz           & 800\,MHz (nominal)       & Teensy 4.1 (+OC) & max. $f_\text{clk}$\\
    On-Chip Flash          & 8\,MB                    & Teensy 4.1       & max. $M_F$ \\
    On-Chip SRAM           & 2\,MB                    & Teensy 4.1       & max. $M_{RAM}$ \\
    Ext. PSRAM             & 8\,MB                    & ESP32-S3         & max. Adressraum\\
    Wi-Fi                  & 802.11\,ax (Wi-Fi\,6)    & ESP32 (Basis)    & best. Standard \\
    Bluetooth              & 5.3 Thread UWB       & nRF52840+        & max. BT-Stack \\
    ADC-Auflösung          & 16-bit, 4\,MSPS          & extern (best.)   & max. $b_{ADC}$\\
    DAC                    & 12-bit, 2-Kanal           & STM32 / ESP32    & analog out \\
    GPIO                   & 80 konfigurierb. Pins    & Teensy 4.1 Basis & max. $n_{GPIO}$\\
    UART/SPI/I2C           & je 4 unabh. Instanzen    & Teensy 4.1       & max. Peripherie \\
    DMA                    & 64 Kanäle, 32-bit         & Teensy 4.1       & max. Durchsatz \\
    USB                    & USB\,3.2 Gen\,1   & kombiniert       & max. USB\\
    EtherCAT-Slave         & integriert (ESC)          & EHAE/OMCU        & Industrie-IIoT \\
    Energieeffiz. aktiv & $< 1\,\text{mW/DMIPS}$   & nRF52840         & min. $P/\text{DMIPS}$\\
    Schlafmodus            & $< 1\,\mu\text{A}$        & MSP430           & min. $P_\text{sleep}$\\
    Betriebsspannung       & 1,8 -- 3,6\,V             & nRF52840/MSP430  & breiter Bereich \\
    Preis & $< 8\,\text{USD}$         & RP2040-Basis     & max. Marktzug.\\
    \bottomrule
  \end{tabular}
\end{table}

\subsection{Architektur-Haupttheoreme}

\begin{theorem}[Existenz und Eindeutigkeit des OMCU-Vektors]
\label{thm:omcu_existenz}
  Sei $\mathcal{M}$ eine endliche nichtleere Menge von Mikrocontrollern
  und sei~$\hat{q}_i(\mu) \in [0,1]$ für alle~$i,\mu$.
  Dann existiert ein eindeutiger Vektor
  $\vec{v}^* \in [0,1]^9$ mit
  \[
    v^*_i \;=\; \max_{\mu \in \mathcal{M}} \hat{q}_i(\mu)
    \quad \forall\, i \in \{1,\ldots,9\},
  \]
  der den utopischen Punkt \omcu{} beschreibt.
\end{theorem}

\begin{proof}
  Da $\mathcal{M}$ endlich ist, existiert für jedes~$i$ ein
  $\mu^*_i \in \mathcal{M}$ mit
  $\hat{q}_i(\mu^*_i) = \max_\mu \hat{q}_i(\mu)$
  (Maximum über endliche Menge).
  Die Eindeutigkeit folgt daraus, dass das Maximum über eine endliche Menge
  eindeutig ist (bei Gleichstand wird dasselbe Supremum angenommen).
 
\end{proof}

\begin{theorem}[Pareto-Superiorität des OMCU]
\label{thm:omcu_pareto}
  Für alle $\mu \in \mathcal{M}$ gilt $\omcu{} \succ \mu$.
\end{theorem}

\begin{proof}
  Sei $\mu \in \mathcal{M}$ beliebig.
  Aus \cref{thm:omcu_existenz}: $\hat{q}_i(\omcu{}) \ge \hat{q}_i(\mu)$
  für alle~$i$.

  Angenommen, $\hat{q}_i(\omcu{}) = \hat{q}_i(\mu)$ für \emph{alle}~$i$.
  Dann müsste $\mu$ in allen neun Dimensionen simultan das Maximum aller
  $\mathcal{M}$-Elemente erreichen.
  Sei~$P_\text{typ}$ der Leistungsparameter:
  $\hat{P}^{-}(\text{MSP430}) = 1$ (vgl.\ \cref{thm:msp430_energie}),
  aber $\hat{f}_\text{clk}(\text{MSP430}) \approx 0{,}01$ (16\,MHz von 600\,MHz).
  Ebenso gilt für den Teensy:
  $\hat{f}_\text{clk}(\text{Teensy}) = 1$, aber
  $\hat{P}^{-}(\text{Teensy}) = 0$ (1500\,mW ist Maximum).
  Es gibt also kein einzelnes $\mu \in \mathcal{M}$, das in allen neun
  Dimensionen gleichzeitig das Optimum hält.
  Daher existiert stets ein~$j$ mit $\hat{q}_j(\omcu{}) > \hat{q}_j(\mu)$.
 
\end{proof}

\begin{korollar}[Maximale Bewertung des OMCU]
\label{kor:omcu_max_phi}
  Für jeden Gewichtungsvektor $\vec{w}$ mit $w_i \ge 0$ und $\sum w_i = 1$:
  \[
    \Phi_{\vec{w}}(\omcu{}) = 1.
  \]
\end{korollar}

\begin{proof}
  Da $\hat{q}_i(\omcu{}) = 1$ für alle~$i$:
  $\Phi_{\vec{w}}(\omcu{}) = \sum_i w_i \cdot 1 = \sum_i w_i = 1$.
 
\end{proof}

\subsection{Radar-Profil und architektonische Visualisierung}

\begin{figure}[H]
  \centering
  \includegraphics[width=0.72\textwidth]{fig6_radar.pdf}
  \caption{Radar-Profil: OMCU (rot, ausgefüllt) gegen ESP32, Teensy\,4.1 und MSP430.
           Der OMCU erreicht in allen acht Dimensionen den Maximalwert\,1.}
  \label{fig:radar}
\end{figure}

\begin{figure}[H]
  \centering
  \includegraphics[width=\textwidth]{fig7_arch.pdf}
  \caption{Blockdiagramm der OMCU-Systemarchitektur.}
  \label{fig:arch}
\end{figure}

\subsection{Formale Komplexitätsbetrachtung}

\begin{proposition}[Berechenbarkeit der OMCU-Synthese]
\label{prop:komplexitaet}
  Die Berechnung des utopischen Vektors~$\vec{v}^*$ hat
  Zeitkomplexität~$O(|\mathcal{M}| \cdot |\mathcal{P}|)$.
\end{proposition}

\begin{proof}
  Für jeden der $|\mathcal{P}| = 9$ Parameter wird einmal das Maximum über
  $|\mathcal{M}| = 12$ MCUs berechnet. Das Maximum über~$n$ Elemente benötigt
  $O(n)$ Vergleiche. Gesamtkomplexität:
  $O(|\mathcal{M}| \cdot |\mathcal{P}|) = O(12 \cdot 9) = O(1)$
  (asymptotisch konstant für feste Eingabegröße, polynomiell in der
  allgemeinen Parameterisierung).
\end{proof}

\begin{proposition}[NP-Vollständigkeit der optimalen Implementierungsauswahl]
\label{prop:np_implementierung}
  Das Entscheidungsproblem
  ``Existiert eine Teilmenge $S \subseteq \mathcal{M}$ mit $|S| \le k$,
  die zusammen die OMCU-Parameter $\vec{v}^*$ vollständig abdeckt?''
  ist NP-vollständig.
\end{proposition}

\begin{proof}[Beweisskizze]
  Dies ist eine Instanz des \emph{Set Cover}-Problems:
  Jeder Mikrocontroller~$\mu_i$ deckt die Menge der Dimensionen ab, in
  denen er den Maximalwert~$1$ hält.
  Set Cover ist NP-vollständig (Karp 1972). Durch Reduktion folgt die
  NP-Vollständigkeit des obigen Problems.
 
\end{proof}

% ══════════════════════════════════════════════════════════════════════════════
\section{Referenzimplementierungskonzept}
\label{sec:implementierung}

\subsection{Hardwareplattform}

Die \omcu{}-Spezifikation aus \cref{sec:omcu} lässt sich durch eine
heterogene Multi-Chip-Architektur realisieren:

\begin{enumerate}[label=\textbf{(\Alph*)}]
  \item \textbf{Recheneinheit:} NXP i.MX\,RT1176 (ARM Cortex-M7 + M4,
        1\,GHz/400\,MHz, 2\,MB SRAM) -- maximale Rechenleistung.
  \item \textbf{Konnektivität:} Infineon CYW43439 (Wi-Fi\,4 + BT\,5.0) oder
        NXP IW612 (Wi-Fi\,6 + BT\,5.3 + UWB).
  \item \textbf{Energiemanagement:} TI BQ25895 PMIC mit MSP430-basiertem
        Schlafmodus-Controller (ultra-low-power Wächter).
  \item \textbf{Speicher:} IS66WVS4M8ALL (32\,MB opsRAM) +
        W25Q128 (16\,MB NOR Flash).
  \item \textbf{Industrieschnittstelle:} BECKHOFF ET1100 EtherCAT Slave
        Controller -- Industrie\,4.0-Kompatibilität.
\end{enumerate}

\subsection{Software-Stack}

\begin{lstlisting}[language=C, caption={OMCU Boot-Sequenz (konzeptionell)}]
/* OMCU Boot-Sequenz: Cortex-M7 Primary Core */
void omcu_boot_sequence(void) {
    power_management_init();      /* MSP430-artiger Ultra-Low-Power Modus */
    clock_system_init(800000000); /* 800 MHz PLL-Konfiguration            */
    cache_enable(L1_ICACHE | L1_DCACHE | TCM_512KB);
    dma_init(64);                 /* 64-Kanal DMA-Controller              */
    wifi6_init(SSID_MANAGER);     /* 802.11ax Association                 */
    ble53_init(BLE_PERIPHERAL);   /* BLE 5.3 + Thread Stack               */
    ethercat_slave_init(ESC_ET1100); /* EtherCAT Slave Controller         */
    cortex_m4_launch(AUDIO_DSP_FIRMWARE); /* Sekundaerkern starten        */
    rtos_start(FREERTOS_CONFIG);  /* FreeRTOS Scheduler                   */
}
\end{lstlisting}

\subsection{Verifikationsansatz via Subgraph Algorithmus}

\begin{definition}[MCU-Eigenschaftsgraph]
\label{def:eigenschaftsgraph}
  Sei $G_\mu = (V_\mu, E_\mu)$ ein Graph, wobei
  $V_\mu = \mathcal{P}$ die Menge der Parameter und
  $E_\mu \subseteq \mathcal{P} \times \mathcal{P}$ die
  funktionalen Abhängigkeiten zwischen Parametern beschreibt
  (z.\,B. $f_\text{clk} \to P_\text{typ}$).
  $G_{\omcu{}}$ ist der vollständige Eigenschaftsgraph des optimalen
  Mikrocontrollers.
\end{definition}

\begin{theorem}[Subgraph-Einbettbarkeit der MCU-Graphen]
\label{thm:subgraph}
  Für jeden Mikrocontroller $\mu \in \mathcal{M}$ ist $G_\mu$
  isomorph zu einem Subgraphen von~$G_{\omcu{}}$:
  $G_\mu \hookrightarrow G_{\omcu{}}$.
\end{theorem}

\begin{proof}
  Da $G_{\omcu{}}$ über denselben Knotensatz~$\mathcal{P}$ und alle
  möglichen Kanten (vollständiger Graph über~$\mathcal{P}$) verfügt,
  ist jeder Graph $G_\mu \subseteq K_{|\mathcal{P}|}$ ein Subgraph von
  $G_{\omcu{}} = K_{|\mathcal{P}|}$, und die Identitätsabbildung
  $\varphi\colon V_\mu \to V_{\omcu{}}$ liefert den Isomorphismus.
 
\end{proof}

Dieser Satz stellt die Verbindung zum Subgraph Algorithmus her
\cite{epp_subgraph2026}:
Die Verifikation, ob eine reale Implementierung die OMCU-Spezifikation
erfüllt, kann als Subgraphiso\-morphie-Problem zwischen dem
Implementierungsgraph und~$G_{\omcu{}}$ formuliert und in
$O(n^3)$ gelöst werden.

% ══════════════════════════════════════════════════════════════════════════════
\section{Kritische Würdigung und Fazit}
\label{sec:fazit}

\subsection{Zusammenfassung der Hauptergebnisse}

Diese Arbeit hat folgende Beiträge geleistet:

\begin{enumerate}[label=(\roman*)]
  \item \textbf{Formales Bewertungsmodell:}
        Die gewichtete Bewertungsfunktion~$\Phi_{\vec{w}}$ und die
        $\ell^2$-Optimalitätsdistanz~$d_{\vec{w}}$ wurden eingeführt,
        ihre Wohldefinierheit und Äquivalenz für die Rangbildung bewiesen
        (\cref{lem:wohldef,thm:aquivalenz}).

  \item \textbf{Vergleichende Analyse:}
        Zwölf MCU-Plattformen wurden anhand von neun technischen
        Parametern normiert und bewertet (\cref{tab:scores}).
        Die plattformspezifischen Stärken (MSP430: Energie,
        Teensy: Performance, ESP32: Konnektivität) wurden formal
        abgeleitet (\cref{thm:msp430_energie,thm:teensy_superior,lem:esp32_pareto}).

  \item \textbf{OMCU-Synthese:}
        Der synthetische Optimale Mikrocontroller wurde als utopischer
        Punkt der Pareto-Front konstruiert (\cref{def:omcu}).
        Seine Pareto-Superiorität gegenüber allen $\mu \in \mathcal{M}$
        und seine maximale Bewertungsfunktion~$\Phi = 1$ wurden bewiesen
        (\cref{thm:omcu_pareto,kor:omcu_max_phi}).

  \item \textbf{Referenzarchitektur:}
        Eine konkrete Hardwareplattform (NXP i.MX\,RT1176 + IW612 + MSP430 PMU)
        wurde skizziert, die die OMCU-Parameter angenähert realisiert.
\end{enumerate}

\subsection{Limitierungen}

Das vorgestellte Modell abstrahiert von:
\begin{itemize}
  \item \textbf{Ökosystem und Toolchain:}
        Bibliotheksreife, Debugging-Support und
        Community-Größe sind schwer zu normieren.
  \item \textbf{Thermisches Design:}
        Hochfrequente Cortex-M7-Cores erfordern Wärmemanagement,
        das in eingebetteten Umgebungen limitierend wirken kann.
  \item \textbf{Lieferkettenstabilität:}
        Chipverfügbarkeit (vgl.\ Halbleiterkrise 2021--2023) ist kein
        technischer aber ein praxisrelevanter Parameter.
\end{itemize}

\subsection{Schlussfolgerung}

Die Synthese des \omcu{} ist kein Plädoyer für ein universelles
Allzweckwerkzeug, sondern ein formales Werkzeug zur systematischen
Anforderungsanalyse: Jede reale MCU-Wahl nähert sich dem \omcu{}
im für die Anwendung relevanten Teilraum von~$\mathcal{P}$.
Die Verbindung zum Subgraph Algorithmus (\cref{thm:subgraph}) öffnet
einen direkten Weg zur formalen Spezifikationsverifikation eingebetteter
Systeme -- ein Bereich, der in zukünftigen Arbeiten weiterverfolgt wird.

\begin{thebibliography}{99}

\bibitem{atmel_atmega328p}
Microchip Technology Inc.,
\textit{ATmega328P Datasheet -- 8-bit AVR Microcontroller with 32K Bytes In-System Programmable Flash},
DS40002061B, 2020.

\bibitem{espressif_esp32}
Espressif Systems,
\textit{ESP32 Technical Reference Manual},
Version 5.3, 2023.

\bibitem{raspberrypi_rp2040}
Raspberry Pi Ltd.,
\textit{RP2040 Datasheet},
Version 2.1, 2023.

\bibitem{st_stm32f103}
STMicroelectronics,
\textit{STM32F103x8/B Reference Manual -- ARM-based 32-bit MCU},
RM0008, Rev.~21, 2021.

\bibitem{pjrc_teensy41}
PJRC Electronics,
\textit{Teensy 4.1 -- i.MX RT1062 Processor Overview},
\url{https://www.pjrc.com/teensy/teensy41.html}, 2023.

\bibitem{ti_msp430}
Texas Instruments,
\textit{MSP430x2xx Family User's Guide},
SLAU144K, 2013.

\bibitem{nordic_nrf52840}
Nordic Semiconductor,
\textit{nRF52840 Product Specification},
v1.7, 2022.

\bibitem{microchip_pic16}
Microchip Technology Inc.,
\textit{PIC16F87XA Datasheet},
DS39582C, 2003.

\bibitem{atmel_attiny85}
Microchip Technology Inc.,
\textit{ATtiny25/45/85 Datasheet},
DS40002088A, 2020.

\bibitem{arm_cortexm7}
ARM Limited,
\textit{Cortex-M7 Processor Technical Reference Manual},
Revision r1p2, ARM DDI0489F, 2019.

\bibitem{pareto_vilfredo}
V. Pareto,
\textit{Manuale di Economia Politica con una Introduzione alla Scienza Sociale},
Società Editrice Libraria, Mailand, 1906.

\bibitem{karp1972}
R.~M. Karp,
``Reducibility Among Combinatorial Problems,''
in \textit{Complexity of Computer Computations}, R.~E. Miller, J.~W. Thatcher, Eds.,
Plenum Press, New York, 1972, pp.~85--103.

\bibitem{epp_subgraph2026}
S. Epp,
\textit{Der Subgraph Algorithmus}, 2026, Verfügbar auf \url{https://github.com/hjstephan86/subgraph}.

\bibitem{cormen2009}
T.~H. Cormen, C.~E. Leiserson, R.~L. Rivest, C.~Stein,
\textit{Introduction to Algorithms}, 3rd~ed.,
MIT Press, Cambridge, MA, 2009.

\bibitem{patterson_hennessy}
D.~A. Patterson, J.~L. Hennessy,
\textit{Computer Organization and Design: ARM Edition},
Morgan Kaufmann, Waltham, MA, 2016.

\end{thebibliography}

\end{document}

